\section{Related work}\label{sec:related}

Our literature search was targeted, not exhaustive. Each comparison is with the statement we checked in the cited source.

\subsection{Monotone arithmetic and tropical circuits}

Valiant~\cite{Valiant80,Valiant79} proved an exponential monotone lower bound for perfect matchings of a triangular grid, gave a polynomial-size general circuit for the same polynomial via Pfaffians, and left rectangular grids unresolved. Jerrum and Snir~\cite{JerrumSnir82} proved exponential monotone lower bounds for the dense permanent and for spanning trees of complete directed graphs, and related $(\min,+)$ and monotone arithmetic complexity. Jukna~\cite{Jukna16Tropical} shows that pure tropical complexity of a homogeneous $0/1$ support family equals its Minkowski (support) complexity, which also lower-bounds monotone arithmetic circuits with arbitrary positive coefficients. His book~\cite{JuknaTropical23} surveys these transfers. Gashkov and Sergeev~\cite{GashkovSergeev12} give a support-based method for monotone lower bounds on explicit polynomials. Chattopadhyay, Datta and Mukhopadhyay~\cite{ChattopadhyayDattaMukhopadhyay21} prove monotone lower bounds via communication complexity, including $\varepsilon$-sensitive versions. Their follow-up with Ghosal~\cite{CDGM22} gives $2^{\Omega(n)}$ monotone lower bounds for directed spanning-tree polynomials of constant-degree expanders, with separate variables for opposite arcs. The surveys of Shpilka and Yehudayoff~\cite{ShpilkaYehudayoff10} and the projection results of Grochow~\cite{Grochow17} concern other questions. Grochow's matching-polytope obstruction is not available for bipartite graphs, whose perfect matching polytope has a small description. Cavalar, Fabris, Mukhopadhyay, Srinivasan and Yehudayoff~\cite{CavalarFabrisEtAl25} give polynomials with positive coefficients that have polynomial-size depth-three circuits but need monotone circuits of size $2^{\Omega(n)}$. This is the strongest possible exact separation and improves on Valiant's. Their basic graph witness is built from $(1-K_G(x))^2$, where $K_G$ counts the edges induced by the selected vertices (Section~\ref{sec:related-pc} explains why such witnesses are not robust after normalization). In a subsection on planar graphs they also note that it is unknown whether deciding perfect matching on planar graphs is hard for monotone \emph{Boolean} circuits. For general graphs, Cavalar, G\"o\"os, Riazanov, Sofronova and Sokolov~\cite{CavalarGoosEtAl25} prove a $2^{n^{\Omega(1)}}$ monotone Boolean lower bound for perfect matching. These Boolean questions concern the existence function of an input graph, not the exact-set indicator studied here.

Curticapean and Wang~\cite{CurticapeanWang26} show that the determinant of an $m\times m$ matrix equals the perfect matching polynomial of a $3m\times5m$ grid with weights in $\{-1,0,1\}\cup\{1-x_{ij}\}$. The construction relates grid matchings to determinants in general arithmetic complexity. It uses negative weights, so it gives no monotone lower bound and is consistent with Theorem~D.

We did not find a later resolution of the square-grid case of Valiant's question. Neither Jukna's papers nor the surveys we checked call it folklore.

\subsection{Knowledge compilation}

Bova, Capelli, Mengel and Slivovsky~\cite{BovaCapelliMengelSlivovsky16} proved DNNF lower bounds with rectangle covers and multi-partition communication complexity~\cite{DurisEtAl04}. Amarilli, Capelli, Monet and Senellart~\cite{AmarilliCapelliMonetSenellart20} give exponential-in-treewidth DNNF lower bounds for monotone CNFs of bounded arity and degree. Ledaguenel, Hudelot and Khouadjia~\cite{LedaguenelEtAl25} apply this to the constraint ``the selected edges form a matching'' on $k\times k$ grids. That is the set of \emph{all} matchings, not the perfect ones, and a lower bound does not pass to a sublayer. De Colnet and Mengel~\cite{DeColnetMengelTseitin21} prove DNNF lower bounds for satisfiable Tseitin formulas, from which a lower bound for perfect matchings of a planar nonbipartite gadget family follows. For the square grid, an exact-support bound follows separately from Theorem~A and Lemma~\ref{lem:balanced-decomp}. If a nonnegative smooth decomposable PC with leaf arity at most $2n/3$ has support exactly $\PM_L$, each nonzero balanced product in the decomposition lies in one profile rectangle: fixing an assignment in either factor forces all assignments in the other to cover the complementary vertex set. These rectangles cover $\PM_L$, so a union bound gives $\abs C\ge1/\alpha_L$. We have not established priority for that exact specialization.

\subsection{Probabilistic circuits}\label{sec:related-pc}

\paragraph{Exact separations.} Martens and Medabalimi~\cite{MartensMedabalimi14} prove the balanced-product theorem for decomposable sum-product networks (their Theorem~38) and an exponential lower bound for uniform spanning trees of the complete graph, exactly and for approximation ``arbitrarily well'' by bounded-size networks (Theorem~37). The latter is not a fixed-TV statement. Zhang, Holtzen and Van den Broeck~\cite{ZhangHoltzenVdB20} relate PCs to determinantal point processes. Zhang, Juba and Van den Broeck~\cite{ZhangJubaVdB21} introduce PGCs and separate them exactly from decomposable PCs using spanning trees. Harviainen, Ramaswamy and Koivisto~\cite{HarviainenEtAl23} study inference in PGCs. Agarwal and Bl\"aser~\cite{AgarwalBlaser24} and Broadrick, Zhang and Van den Broeck~\cite{BroadrickZhangVdB24} relate PGCs to signed circuits (Section~\ref{sec:signed}). Loconte et al.~\cite{LoconteEtAl24}, Loconte, Mengel and Vergari~\cite{LoconteMengelVergari24} and Wang and Van den Broeck~\cite{WangVdB25} give exact separations between monotone, squared and sum-of-squares circuits. These are mostly \emph{structured}, and their witnesses are mostly unique-disjointness functions. Zhao et al.~\cite{ZhaoEtAl26} prove separation-rank bounds for structured-decomposable PCs with a fixed vtree. Zellinger et al.~\cite{ZellingerEtAl26} study approximate inference for subtractive mixtures.

\paragraph{Robust separations.} Leland and Choi~\cite{LelandChoi25} prove that \emph{deterministic} decomposable PCs need size $2^{\Omega(n)}$ to approximate a certain distribution within a constant TV distance, while a polynomial-size nondeterministic decomposable PC represents it exactly. Their lower bound requires determinism. Theorem~B does not, and it uses a different mechanism (an overlap bound rather than pruning to Boolean approximations).

\paragraph{Exact separations can vanish after normalization.} A lower bound for exact representation of a function says little about approximating the normalized distribution. Here is the standard unique-disjointness witness. Let $K(x,y)=\sum_{i=1}^mx_iy_i$ on $\{0,1\}^{2m}$, $f=(1-K)^2$ and $g=K^2+1=f+2K$. Let $P_f$ and $P_g$ be the distributions proportional to $f$ and $g$. Expanding $K^2$ gives a smooth decomposable PC for $g$ with $O(m^2)$ nonnegative terms, each a product of at most four variables. Under the uniform measure $\E[K]=m/4$, and $P_g$ is the mixture of $P_f$ and $P_{2K}$ with weight $1-\E[f]/\E[g]$ on the second component. So
\[
  \TV(P_f,P_g)\le1-\frac{\E f}{\E g}=\frac{m/2}{\E[K^2]+1}\le\frac{m/2}{m^2/16+1}\le\frac8m .
\]
The same calculation applies to graph-polynomial witnesses of the form $(1-K_G)^2$, such as the shallow construction of Cavalar et al.~\cite{CavalarFabrisEtAl25}. So separations driven by cancellation at a small set of points do not, by themselves, give constant-TV lower bounds. In the dimer case the hard set is the whole support.

Robust versions of unique disjointness need the same care. Braun and Pokutta~\cite{BraunPokutta16} show that the nonnegative rank stays exponential under a large class of changes: every nonnegative matrix indexed by pairs of subsets of $[N]$, with entries $r$ on disjoint pairs, $r-1$ on pairs meeting in one element and arbitrary nonnegative entries elsewhere, has nonnegative rank at least $2^{N/(8r)}$ (Theorem~4.4 in ECCC TR13-056, revision~5). This is robust entrywise. It does not by itself give a constant-TV bound for a normalized distribution. Under the uniform measure on pairs, only a $(3/4)^N+\tfrac N4(3/4)^{N-1}$ fraction of pairs meet in at most one element, so the constrained entries can carry a vanishing share of the mass. We do not use this theorem in our proofs.

\subsection{A stronger generic separation from codes}\label{sec:codes}

De Colnet and Mengel~\cite[Sec.~5]{DeColnetMengel20} construct, using the multi-partition rank bounds of~\cite{DurisEtAl04}, binary matrices $H\in\mathbb F_2^{k\times n}$ with $k=n/100$ such that every submatrix with at least $n/3$ columns has rank at least $k-1$. Their focus is approximate knowledge compilation, so they do not state the following corollary. Let $P_{\mathcal C}$ be uniform on the code $\mathcal C=\ker H$. For a balanced cut $A\sqcup B$ of the columns, the syndrome $z=H_Ax_A=H_Bx_B$ plays the role of the profile. Each syndrome class is a product set, and all of them have probability $2^{-(\rank H_A+\rank H_B-\rank H)}\le2^{-(k-2)}$. Lemma~\ref{lem:overlap}, with profiles replaced by syndromes, and Lemma~\ref{lem:balanced-decomp} show that every nonnegative smooth decomposable PC with leaf arity at most $2n/3$ within TV $\eps$ of $P_{\mathcal C}$ has size at least $(1-\eps)^22^{k-2}$. This is $\exp(\Omega(n))$ for every fixed error below one. Meanwhile, Gaussian elimination computes every evidence probability and conditional of $P_{\mathcal C}$ exactly. This gives a stronger, $\exp(\Omega(n))$, constant-TV separation between nonnegative decomposable PCs and a tractable model. It uses a random, nonlocal construction. Our results concern a specific natural distribution and give $\exp(\Theta(\sqrt n))$, which is tight for it.

\subsection{Dimers, kernels and entanglement}\label{sec:related-physics}

The dimer model on the square grid was solved by Kasteleyn~\cite{Kasteleyn61} and Temperley and Fisher~\cite{TemperleyFisher61,Fisher61}. Inverse Kasteleyn asymptotics in infinite volume and for Temperleyan domains are due to Kenyon~\cite{Kenyon97,Kenyon00}; see also his lectures~\cite{KenyonLectures}. Thurston's height functions~\cite{Thurston90} give another view of domino tilings. Stéphan, Misguich and Pasquier~\cite{StephanMisguichPasquier12} relate R\'enyi entropies of RK states to classical boundary probabilities. They study semi-infinite cylinders and the Kitaev--Preskill disk geometry, and they interpolate between triangular and square lattices with a bond fugacity. Their straight-boundary and special-geometry results do not cover arbitrary link bipartitions. Starr and Williams~\cite{StarrWilliams18} bound an emptiness formation probability for aligned blocks using reflection positivity. Corollary~\ref{cor:entanglement} is a min-entropy bound over \emph{all} balanced link bipartitions of the free square, including disconnected and interlaced ones. We did not find this statement in the literature we checked. Single-copy entanglement was introduced in~\cite{EisertCramer05}.

The ingredients of the proof of Theorem~A are individually standard: tilted Kasteleyn determinants, the Carleman-type bound and inverse-kernel estimates. We did not find them combined with row-spread forbidden sets to bound arbitrary-cut profiles.
